\chapter{Last Scattering and Baryogenesis}

\section{Looking Outward Means Looking Backward}

At different distances we see the universe at different times. Nearby, the look-back time is not very important. The universe a year ago, or a hundred years ago, was not very different from the universe now. But far enough out, we see the universe younger. A younger universe is denser, and a denser universe is hotter.

At some sufficiently large distance, the universe was about as hot as the surface of the sun. It does not look that way to us because the photons have been redshifted on the way here. A redshift factor of roughly a thousand means an observed temperature lower by the same factor:
\[
T_{\mathrm{obs}} \sim \frac{T_{\mathrm{emit}}}{1+z},
\qquad
1+z \sim 10^3.
\]\footnote{This display is a compact editorial transcription of the verbal redshift argument at 00:02:20--00:02:39.}

That is where the universe became opaque. The universe was mainly hydrogen, with a bit of helium. Above a temperature of a few thousand degrees, hydrogen is ionized. So beyond that surface there are not neutral atoms but ions: mainly protons and electrons moving freely. The photons are also more energetic there, basically X-rays. If an atom forms accidentally, it is quickly reionized again by collisions or by a photon.

An ionized gas is an electrical conductor, and a conductor does not let radiation pass through cleanly. Charges respond to electric fields, slosh back and forth, and re-emit radiation. So beyond that distance the universe is opaque to us. The surface of last scattering is like the visible surface of the sun: not a material wall, but a place where the medium rather suddenly goes from opaque to transparent.

Most of the particles there are photons, by a factor of about \(10^8\). The same is true today, although today the photons are so low in energy that they are not an important part of the present energy density. Once the temperature drops below the ionization temperature of hydrogen, the universe becomes clear and we can see through it.

Direct sight ends there. The next question is how much farther back known atomic, nuclear, and particle physics can still take us.

\section{Going Farther Back}

Just inside last scattering, in addition to photons, the main particles are protons and electrons, with some helium nuclei and a few deuterons. There are practically no antiprotons and practically no positrons. But farther back, the temperature eventually becomes high enough that two photons can collide and produce an electron and a positron:
\[
\gamma + \gamma \leftrightarrow e^- + e^+.
\]

The basic threshold is that the photons must at least supply the rest energy of the pair:
\[
E \ge 2m_e c^2.
\]
Below that threshold ordinary light does not do the job. Above it, electron--positron pairs can be created, and the inverse reaction can also happen. A hot pot of material at that temperature would contain particle reactions playing the role that chemical reactions play in ordinary matter.

Just above threshold there are still more photons than electrons and positrons. Electron rest mass is expensive. But as the temperature rises far above threshold, the electron rest energy becomes less important than the kinetic energy, and the densities become comparable. Then electrons, positrons, and photons all have roughly the same density and are colliding around like mad.

\begin{figure}[ht]
\centering
\includegraphics[width=0.62\textwidth]{lecture_06_figure_02.png}
\caption{Pair-production threshold sketch comparing the electron and proton mass scales.}
\lectureframe{00:18:47}
\end{figure}

The next question is what happens to protons. They are already present, and by then they are very dense, but there is not yet enough thermal energy to create proton--antiproton pairs in abundance. The reason is simple:
\[
m_p \approx 2000\,m_e,
\qquad
E \ge 2m_p c^2
\]
is a much higher threshold.\editorialnote{In the frame the lower mass label is less sharp than the upper \(2m_e c^2\), but the spoken comparison is explicitly between the electron and proton thresholds.}

This is already much hotter than the interior of a star, where antiparticles are not present in any substantial abundance.

But even that is only a provisional description. By the time the temperature is high enough to make protons and antiprotons freely, it is also high enough to smash protons into quarks. So the really hot soup is better described as a soup of electrons, positrons, quarks, antiquarks, photons, neutrinos, gluons, and other species. A simple schematic reaction from the board is
\[
e^- + e^+ \to \gamma \to q + \bar q.
\]

\begin{classroomqa}
\lecturetimestamp{00:20:05}
\audiencequestion{Is it mesons or pre-quarks?}
\lecturerresponse{Mesons are made out of quarks. Once the temperature is high enough to shatter a proton, it also shatters the mesons.}
\end{classroomqa}

So far the reasoning has been backward in time. Now turn it around and ask what happens as the hot symmetric soup cools.

\section{Electrical Neutrality and a Symmetric Hot Universe}

At sufficiently high temperature the universe contains huge densities of particles and antiparticles. Electrons and positrons occur in pairs, quarks and antiquarks occur in pairs, and the natural picture is nearly equal abundances of each particle and its antiparticle.

Electrical neutrality needs a little care. The universe is electrically neutral today. A useful mnemonic is to imagine, not our actual universe, but a closed and bounded one, like a three-sphere. Think first of a two-dimensional sphere as a model. Electric field lines cannot begin or end except on charge. A positive charge emits lines of force; a negative charge absorbs them. On a closed sphere those lines cannot simply escape to infinity. They refocus at the opposite side. So, as a mnemonic, a closed and bounded universe cannot carry net charge. Whether our actual universe is literally closed is not the point. The point is that large-scale positive and negative charges appear to cancel extraordinarily well.

Electrical neutrality does \emph{not} mean that the number of protons is the same as the number of antiprotons. It means that the total positive charge equals the total negative charge. Protons, positrons, and every other positively charged species together balance electrons and every other negatively charged species together.

As far back as known physics can be trusted, and perhaps a little beyond, the universe was hot enough to sustain both particles and antiparticles in abundance. Any accelerator experiment one can name had an analogue in the early universe at some stage. A nearly symmetric start is the natural conjecture: equal numbers of electrons and positrons, equal numbers of quarks and antiquarks, and, if protons and antiprotons are used as stand-ins for baryons and antibaryons,
\[
N_p = N_{\bar p}.
\]

That conjecture is attractive because the laws of physics are almost symmetric between particles and antiparticles. A positron behaves very much like an electron, not quite but very nearly. That near symmetry is the problem.

\section{Running the Movie Forward}

Suppose the universe begins in a hot state with equal or nearly equal abundances of particles and antiparticles, and then cools.

While the temperature is high enough, protons and antiprotons can be created from energetic collisions, and they can annihilate back into photons. But eventually the temperature falls below the proton threshold. Then proton--antiproton pairs can still annihilate into photons, but after a little expansion those photons have lost enough energy that they can no longer turn back into proton--antiproton pairs. This is the one-way street. The same thing happens later for electrons and positrons, at a much lower temperature than the proton threshold but still far above stellar temperatures.

Because the expansion is slow compared with microscopic reaction rates, the universe is mostly close to thermal equilibrium while this is happening. This is what is meant here by adiabatic expansion. In that setting the more accurate statement is that the entropy changes very little. For a hot particle soup, entropy is essentially the number of particles, up to factors of order one. So the protons, antiprotons, electrons, and positrons that annihilate are roughly replaced by photons.

If the numbers of protons and antiprotons were exactly equal, they would simply annihilate away. The same would happen to electrons and positrons. One would be left with photons. But if there were even a tiny excess of protons over antiprotons, the matched pairs would annihilate and the unmatched excess would remain. A concrete example makes the point: if there were one hundred protons and fifty antiprotons, the fifty antiprotons would find fifty protons, annihilate with them, and fifty protons would remain. Charge neutrality then requires the same sort of leftover excess for electrons over positrons.

This lets one read the early asymmetry from the present photon count. Today the number of photons is about \(10^9\) to \(10^{10}\) times the number of protons or electrons:
\[
\frac{N_\gamma}{N_p} \sim 10^9\text{--}10^{10}.
\]
So in the early hot soup there was only about one extra proton for every billion proton--antiproton pairs. That tiny imbalance is the whole reason atoms still exist.

\begin{classroomqa}
\lecturetimestamp{00:42:25}
\audiencequestion{Why not think of the excess statistically, as something created during cooling with a tiny probability?}
\lecturerresponse{To get an extra proton from an initially symmetric state, one would have to create a proton without also creating an antiproton. One must also create negative charge with it, say an electron. The known standard model contains no such low-energy process. At those tested energies baryon number is conserved, so ordinary cooling does not generate the net excess.}
\end{classroomqa}

So electric charge is not the obstacle. A schematic example would be
\[
\gamma + \gamma \to p + e^-.
\]\footnote{This display is an editorial transcription of the verbal example at 00:44:16--00:46:35. Its purpose is to isolate the point that electric charge could be conserved while baryon number is violated.}
This would be electrically neutral overall. But it violates another empirical rule. In the tested standard model one cannot create a proton without also creating an antiproton. For present purposes the relevant conserved quantity is baryon number,
\[
B=(N_p+N_n)-(N_{\bar p}+N_{\bar n}).
\]\footnote{This compact formula formalizes the verbal definition given at 00:45:18--00:45:40.}
At known energies that quantity is conserved, so the asymmetry cannot be manufactured during the late, well-tested part of the cooling history.

\section{Baryogenesis}

This is the subject called baryogenesis. Suppose the universe really did begin with equal numbers of baryons and antibaryons. What would the laws of physics have to allow in order to tilt that symmetric state into the tiny observed excess?

On the board this was reduced to two symbolic expressions,
\[
N_p=N_{\bar p},
\qquad
N_p-N_{\bar p},
\]
and then to a list of conditions connecting the first to the second.

\begin{figure}[ht]
\centering
\includegraphics[width=0.72\textwidth]{lecture_06_figure_04.png}
\caption{Conditions for generating a baryon asymmetry from an initially symmetric state.}
\lectureframe{01:22:26}
\end{figure}

\editorialnote{In this frame the words in the numbered list are only partly legible. The three items below follow the accompanying spoken discussion.}

The first condition is violation of baryon number. If baryon number is an exact conservation law, a universe that starts with zero net baryon number stays that way.

The second condition is a bias between particles and antiparticles. Baryon-number violation alone is not enough. If nature allows a proton-producing process and its reflected antiproton-producing process with equal probability, the average imbalance is still zero. One needs the rates to differ.

Ordinary statistical fluctuations are also far too small. If the total number of particles is about \(10^{90}\), including photons, then a fluctuation of order \(\sqrt{N}\) gives only about \(10^{45}\), nowhere near enough to explain the roughly \(10^{80}\) protons in the universe.

The third condition is that the universe must at some point have been sufficiently out of equilibrium. Here the argument uses the CPT theorem. If a hot gas has sat in equilibrium forever, a movie run backward looks the same as a movie run forward. In that setting the symmetry between particles and antiparticles is effectively restored by the exact symmetry that combines charge conjugation, parity, and time reversal. To get a real bias frozen in, the actual history has to look different going forward and backward on microscopic time scales. Early enough, the expansion of the universe was fast enough to do that.

So the three conditions are:
\begin{enumerate}
\item violation of baryon number;
\item violation of the symmetry between particles and antiparticles, more precisely CP violation;
\item a stage out of thermal equilibrium.
\end{enumerate}

\begin{classroomqa}
\lecturetimestamp{01:09:08}
\audiencequestion{Matching these to CPT, is condition 2 really CP and condition 3 really the time-direction issue?}
\lecturerresponse{Yes. Condition 2 should really be stated as CP violation. Condition 3 is the fact that the actual history was out of equilibrium, not a separate conservation law and not the same thing as violating baryon number.}
\end{classroomqa}

With these three conditions in place, even a universe that starts with equal numbers of protons and antiprotons can be tilted one way or the other. What remains uncertain is whether nature really violates baryon number, and if so how.

\section{What Is Known, and What Is Not}

The open part of the story is baryon-number violation. We have never directly seen a reaction in nature that changes
\[
(N_p+N_n)-(N_{\bar p}+N_{\bar n}).
\]
That may only mean that the relevant reactions require much higher temperatures than anything directly tested in the laboratory, or that they are fantastically improbable at low energy. But at present this is still the question mark.

That is why proton decay matters so much in the discussion. If baryon number is not exact, then the proton should in principle be unstable. But it is spectacularly stable. The lower bound quoted here is of order
\[
\tau_p \gtrsim 10^{34}\ \text{years}.
\]
This comes from huge experiments using very large volumes of water instrumented with light detectors and buried deep underground. It does not disprove baryon-number violation. It means only that any such violation must be extraordinarily suppressed at low energy.

A separate known-physics remark also enters the discussion. Around \(100\) to \(250\,\mathrm{GeV}\), the weak interactions become much more like the electromagnetic interactions. That may matter in detailed models, but it is not the dramatic step in the present argument.

Dark matter does not change the basic logic much. In most theories it interacts so weakly with ordinary matter that, for this question, it is essentially a spectator.

\begin{classroomqa}
\lecturetimestamp{01:24:49}
\audiencequestion{How does gravity affect all this?}
\lecturerresponse{Mainly in two ways. First, without cosmic expansion there would be no out-of-equilibrium stage. Second, black-hole evaporation gives a strong reason to suspect that baryon number cannot be an exact conservation law forever. Apart from that, ordinary gravitational interactions between particles are too weak to matter much here.}
\end{classroomqa}

There is then a follow-up question about electric charge. Would the same black-hole argument also threaten charge conservation? The answer is no. A charged black hole carries electric flux lines extending outward, and those lines keep track of the charge. Such a black hole does not evaporate all the way away; it approaches an extremal remnant instead. So the black-hole argument threatens baryon number, not electric charge.

Another question asks whether the very dense early quark soup should simply be regarded as a black hole. The answer is no. The point is not a static dense lump. The universe was expanding. The black-hole story is a thought experiment about the fate of baryon number, not a literal description of the early universe as a trapped region.

\begin{classroomqa}
\lecturetimestamp{01:31:14}
\audiencequestion{Would one know how to amend the standard model to include such reactions?}
\lecturerresponse{The trouble is the opposite. It is easy to add too much baryon violation and make the proton decay far too quickly. The real problem is keeping the proton as stable as it is while still allowing the early-universe processes one wants.}
\end{classroomqa}

\section{Entropy and the Late Neutron Questions}

Return now to the bookkeeping. In an adiabatically expanding universe, entropy is approximately conserved. For a hot soup of elementary particles, entropy is essentially the particle number up to factors of order one. That is why the discussion does not need a huge late production of photons out of nowhere. The annihilations mostly reshuffle the content of the soup into radiation.

The older version of the problem had the story upside down. Instead of asking why so few protons survived relative to photons, the question was asked the other way around: why are there so many photons compared with protons? But the sharper question is not why the universe manufactured so much entropy late in the game. It is why the asymmetry between particles and antiparticles was so tiny to begin with: only about one extra proton per billion proton--antiproton pairs.

A free neutron does decay, but that process does \emph{not} violate baryon number:
\[
n \to p + e^- + \bar\nu.
\]
The neutron carries the same baryon number as the proton.

For a free neutron the available decay energy is very small, because the neutron mass is only slightly above the mass of the final state. That is part of why the decay is so slow by elementary-particle standards. When a neutron is bound in a nucleus, binding energy can remove that small amount of available energy altogether, and then the bound neutron can be stable.

\begin{classroomqa}
\lecturetimestamp{01:43:03}
\audiencequestion{Are neutrons simply disappearing from the universe, or are they also created?}
\lecturerresponse{The first answer is that free neutrons are lost through decay. But that is not the whole story. There are real sources of neutron production, including nuclear reactions in stars and the formation of neutron stars. So the total neutron count is not simply a one-way decrease.}
\end{classroomqa}
